我叫马睿骁，目前为浙江理工大学物理学专业硕士研究生，预计于2027年6月毕业。硕士阶段，我主要围绕低维功能材料、多铁耦合及相关量子物性开展第一性原理计算和多尺度模拟研究。

刚开始接触科研时，我更多关注如何完成具体计算、如何获得可靠结果。随着研究不断深入，我逐渐发现，真正吸引我的并不是某一种计算方法本身，而是计算结果背后的物理问题：原子结构为什么会发生变化，极化如何影响电子态和微观相互作用，界面、应变和局域结构又如何改变材料的功能性质。

这些经历使我的研究兴趣逐渐从单一物性的计算扩展到铁电极化、结构相变、界面效应以及不同自由度之间的耦合问题。

硕士阶段的研究主要以理论计算为主，我已经积累了较为系统的第一性原理、科研编程、数值模拟和机器学习基础。博士阶段，我希望继续发挥这些优势，但相比单纯延续纯计算研究，我更希望能够真正参与材料制备和实验表征，理解理论模型中的结构、极化和相变如何在真实样品中出现，并尝试建立理论与实验相互支持的研究方式。

\section{一、专业基础}

硕士期间，我接受了较系统的凝聚态物理和计算材料研究训练。目前较为熟悉 VASP、Wannier90/WannierTools、TB2J、Spirit 等计算工具，能够开展结构优化、电子结构、自旋轨道耦合、铁电极化及翻转路径、界面结构、磁交换作用、磁各向异性、Berry曲率及相关拓扑性质等计算，并能够根据具体问题建立相应模型和开展进一步的数值模拟。

在长期研究过程中，我逐渐形成了一种比较明确的分析思路：

从晶体结构出发，分析电子结构和轨道行为，再进一步追踪微观相互作用，并最终理解材料功能性质的来源。

相比单纯得到一个符合预期的结果，我更加关注这个结果是否可靠。因此，对于 U 值、SOC、k点密度、Wannier拟合窗口、有限尺寸效应等可能影响结论的因素，我通常会主动进行参数测试和对照计算。当不同方法得到不一致结果时，我也会重新检查结构、计算条件和模型假设，而不是直接选择一个看起来更“合理”的结果。

\section{二、编程与科研代码开发}

我能够使用 Python、Shell/Linux 编写科研程序和自动化脚本，用于高性能计算任务管理、批量模型生成、参数扫描、结果提取、统计分析和可视化。

对我而言，编程的价值并不只是减少重复操作，而是让我在现有软件无法直接回答某个问题时，可以进一步理解数据和计算方法本身，并尝试自行建立分析工具。

例如，在 TcIrGe₂Se₆ 的研究中，为了理解其中的长程交换机制，仅依靠能带和投影态密度难以定量比较不同轨道通道的贡献。为此，我基于 Wannier 紧束缚模型自行编写 Python 程序，解析 wannier90\_hr.dat、Wannier中心和晶体结构信息，提取不同原子、不同轨道及不同晶格平移之间的跃迁矩阵元，并进一步分析特定交换路径中的轨道杂化。

这一过程使原本主要依靠图像和经验判断的问题，转化为可以直接计算和比较的数据。

这也是我在科研过程中逐渐形成的习惯：

如果现有方法不能直接回答问题，就进一步理解问题本身，再尝试通过程序把它转化成可以计算和验证的形式。

\section{三、科研经历与创新实践}

硕士期间，我主要围绕二维多铁材料和低维功能材料开展研究。

在二维多铁材料 TcIrGe₂Se₆ 的研究中，我首先分析材料的结构稳定性和铁电极化，在此基础上进一步研究不同极化状态下的电子结构、磁性、谷自由度和Berry曲率等性质。

为了进一步理解其中微观相互作用的来源，我构建 Wannier 紧束缚模型，对 Tc–Se–Ir–Se–Tc 长程超超交换通道进行分析，将局域结构变化、轨道杂化和微观交换作用联系起来。

通过这一工作，我逐渐形成了从

结构变化 → 极化状态 → 电子结构 → 轨道杂化 → 微观相互作用 → 功能性质

逐层分析材料问题的研究方式。

相关工作 “Ferroelectric control of magnetism, valley polarization, and skyrmions in monolayer TcIrGe₂Se₆” 已被 Physical Review B 接收。

除此之外，我还开展了成分替位调控二维多铁材料性质的研究，并将第一性原理得到的微观参数进一步用于有限温度模型分析。目前，我正在推进二维铁电/磁性异质结构研究，主要关注极化翻转、界面构型、层间耦合和应变对材料性质的影响。

随着这些课题不断推进，我的关注重点也逐渐发生变化。相比继续寻找某一种具有特殊性质的新材料，我现在更希望回答：

材料的结构为什么发生变化？这种变化为什么会改变材料性质？这种变化又能否通过极化、应变、界面或者外场主动控制？

这也是我希望博士阶段继续深入的问题。

\section{四、机器学习与新方法探索}

除第一性原理和模型计算之外，我也在主动探索机器学习与材料物理问题的结合。

在交错磁拓扑研究中，我尝试将 Sobol参数采样、拓扑不变量计算、机器学习预测和局部参数搜索结合起来。

面对较大的模型参数空间，如果完全依赖人工逐点扫描，不仅效率较低，也难以及时判断哪些区域真正值得进一步研究。因此，我首先利用数值方法生成参数数据，再通过机器学习模型辅助寻找可能出现拓扑相变的区域，最后重新回到严格的物理计算中进行确认。

通过这一研究，我逐渐形成了一个比较明确的认识：

机器学习不应该替代物理，而应该帮助研究者更高效地发现物理。

未来，我希望进一步学习机器学习原子势等方法。如果能够以第一性原理数据训练可靠的原子势，就有机会把传统DFT能够处理的时间和空间尺度进一步拓展，用于研究大尺度极化翻转、畴壁运动、缺陷作用以及有限温度下的结构演化。

这也能够把我目前已有的第一性原理、编程和机器学习能力进一步结合起来。

\section{五、研究兴趣与博士阶段研究设想}

经过硕士阶段的学习和研究，我目前最希望进一步深入的是铁电/反铁电材料、外延薄膜及异质界面中的结构和极化问题。

首先，我非常关注极化翻转过程中的原子尺度结构变化。

很多理论研究最终比较的是两个稳定极化态，而真实极化翻转往往涉及原子的连续移动、中间亚稳结构甚至缺陷参与。我希望进一步理解不同翻转路径为什么具有不同的能垒，以及应变、界面、缺陷和晶向如何改变整个翻转过程。

其次，我对铁电相、反铁电相以及其他亚稳结构之间的竞争关系很感兴趣。

在真实材料中，不同结构的稳定性往往同时受到外延应变、电场、界面、电荷重构和缺陷等因素影响。我希望结合第一性原理、NEB和分子动力学等方法，对这些过程进行更加系统的研究。

此外，我也希望进一步研究外延薄膜和异质界面。与理想的自由材料相比，真实薄膜中的外延约束、界面结构和局域缺陷往往会带来更加复杂，但也更加接近实际材料的问题。

对于极化畴、畴壁、极化涡旋以及其他复杂极化结构，我同样具有较强兴趣。这些结构涉及晶格、弹性场、静电作用和界面条件之间的竞争，也是连接原子尺度结构和宏观功能响应的重要问题。

\subsection*{关于理论与实验的选择}

虽然我的硕士阶段主要接受的是理论计算训练，而且我也愿意继续发挥这一优势，但博士阶段我更加希望能够真正参与实验工作。

我希望系统学习材料制备、薄膜生长、结构表征和物性测试，亲自了解一个理论预测如何在真实材料中被实现和验证。

如果所在团队的研究模式更适合理论与实验分工，我也非常愿意承担与实验紧密配合的理论计算工作，例如进行材料和结构设计、第一性原理计算以及微观机制分析，为实验提供预测和解释，同时根据实验结果重新调整理论模型。

因此，我对博士阶段最理想的设想并不是完全从理论转向实验，也不是继续只做纯计算，而是：

以现有理论计算和编程能力为基础，尽可能参与实验；同时让自己的计算直接服务于真实的实验问题。

我希望理论能够帮助实验判断“什么值得做、可能出现什么”，实验则能够告诉理论“真实材料究竟发生了什么”。

相比将计算和实验分成两个彼此独立的部分，我更希望围绕同一个科学问题，让二者真正相互推动。

\section{六、博士阶段学习计划}

博士阶段前期，我希望进一步加强凝聚态物理、铁电物理和材料相变等方面的理论基础，并尽快熟悉课题组正在研究的材料体系和实验方法。

在计算方面，我会继续发挥已有的第一性原理和程序开发能力，并进一步学习分子动力学、机器学习原子势以及适用于铁电体系的多尺度模拟方法。

实验能力则是我博士阶段希望重点补充的部分。

我希望有机会系统学习外延薄膜生长、铁电性能测试、PFM、高分辨STEM以及同步辐射结构表征等技术，真正了解从材料制备、结构调控到性能测试的全过程。

博士阶段中期，我希望逐渐围绕具体材料形成相对独立的研究课题，并尝试真正把计算和实验结合起来。

例如，可以先从计算中预测外延应变、界面结构或缺陷可能引起的结构变化，再通过实验进行观察；反过来，如果实验发现理论模型无法直接解释的现象，则重新建立模型和计算方案寻找原因。

我希望最终形成一种自然的研究方式：

计算帮助实验提出问题，实验帮助理论发现新的问题。

博士阶段后期，我希望进一步提高独立凝练科学问题、设计研究方案和推进完整课题的能力，在完成博士论文和相关成果的同时，逐渐形成自己相对稳定的研究方向。

\section{七、职业规划}

对于未来职业选择，我希望保持一定的开放性。

如果博士阶段能够取得较好的科研成果，并有机会进入合适的平台继续开展研究，我愿意在博士毕业后继续进行博士后训练，进一步积累研究成果、合作经验和独立科研能力，并争取进入高校或科研院所。

高校教师和科研人员仍然是我比较向往的发展方向，因为这样的职业能够让我持续接触新的科学问题，并有机会长期开展具有连续性的研究。

但我也清楚，职业规划需要结合博士阶段的实际发展情况。

如果届时自己的科研成果、研究方向或发展机会更适合其他道路，我也会认真考虑高校教学岗位、科研院所或者企业研发等选择。特别是在新材料、半导体、计算材料和人工智能辅助材料研发等方向，企业同样可以提供较好的研究和技术平台。

因此，相比现在就限定自己未来必须从事哪一种职业，我更希望通过博士阶段不断提高自己的科研能力和专业竞争力，使未来拥有更多选择。

如果科研成果、平台和发展机会能够支持我继续走学术道路，我愿意继续向博士后和高校方向发展；如果实际情况更适合其他选择，我也会选择能够发挥自身优势、继续学习并获得更大发展空间的岗位。

无论最终进入高校、科研机构还是企业，我都希望自己的工作能够保留一定的探索性和创造性。

\section{八、寄语——从理论预测走向真实材料}

虽然我的硕士阶段主要从事理论计算，但相比长期停留在纯理论计算中，我其实更加向往能够亲自参与实验。

第一性原理计算为我理解材料提供了非常重要的工具，也构成了我目前最主要的科研基础。但几年的研究也让我逐渐认识到，理论计算始终建立在一定的模型和参数之上，而真实材料中还存在温度、缺陷、应变、界面以及动力学过程等更加复杂的因素。

因此，我越来越希望真正回答一个问题：

理论上预测出的材料和性质，在实验中究竟能不能被做出来、被观察到，并最终被测量到？

如果能够从理论上提出一种可能的材料或结构，通过计算判断其稳定性和物理机制，再亲自参与样品制备和实验表征，最终在实验中观察到与理论预测相对应的现象，我认为这种过程会给我非常直接的科研反馈。

对我而言，这种反馈不仅意味着理论得到了验证，更重要的是说明自己的研究真正和现实材料建立了联系。

理论可以帮助实验减少盲目试错，实验结果又可以反过来检验和修正理论模型。如果实验结果和预测一致，我会享受这种从理论走向现实的过程；如果两者并不一致，我同样愿意继续追问其中的原因，因为很多真正有意义的问题，恰恰隐藏在理想模型与真实材料之间的差异中。

我很喜欢这样一种科研过程：

提出一个想法， 用理论判断它是否可能实现， 再尝试真正把它做出来， 最后通过实验判断最初的想法是否成立。

这也是我目前对博士阶段最向往的研究状态。

同时，我一直希望能够进入研究条件更加丰富、学术交流更加充分的平台。对我而言，更高水平的平台不仅意味着更多设备和资源，更重要的是能够接触更加优秀的研究者、更加前沿的问题，以及过去没有机会使用的方法。

探索未知一直是我非常喜欢的一件事。如果未来的科研环境、个人发展和现实条件允许，我希望能够长期从事科研工作。

博士阶段对我而言，因此不仅是学历上的继续深造，也是一次重新拓展自己研究边界的机会。

我希望自己能够从目前以理论计算为主要基础的硕士生，逐渐成长为能够在理论预测、程序开发、数据分析、材料制备、实验表征和物理解释之间建立联系的研究者，并在不断探索未知的过程中找到适合自己长期发展的方向。
